\section{Further research questions}

\begin{enumerate}
\item \textbf{Direct Grassmann--Cayley certificates.}
Can the Pappus, Pascal, and Brianchon calculations be replaced by short
bracket-algebra derivations whose syntax compiles directly to the native
\texttt{ConciseIPRA} equational calculus, avoiding coordinate normalization
altogether?

\item \textbf{A universal certified-intersection theorem.}
Is there a single list theorem in $\ipra$ saying that endpoint evaluations of
all certified line restrictions telescope around a closed walk, with
Menelaus, Carnot polygon, Ceva-type products, and Hoehn-type identities as
instances?

\item \textbf{Proof-object minimization.}
For each fixed polynomial certificate, what is the smallest proof accepted by
the five-rule or six-rule native calculi in \texttt{ConciseIPRA}?  The vector
identities for Pappus and Pascal may be substantially shorter than determinant
normalization certificates.

\item \textbf{Degenerate closure.}
The division-free identities remain meaningful when some named intersections
coalesce or move to infinity.  Can one formulate a uniform projective theorem
using only proportional triples and incidence, then derive every classical
nondegenerate ratio statement as a corollary?

\item \textbf{Ordered conclusions.}
What is the weakest constructive ordered arithmetic needed to recover the
positive-length versions, betweenness statements, and the convex-pentagon
interpretation of Hoehn from the signed identities?

\item \textbf{Automatic dual theorem generation.}
Can a Lean metaprogram dualize any point--line theorem expressed in the
cross-product kernel and generate both the field-level proof and the coded
$\ipra$ statement?  Pascal-to-Brianchon is the first test case.

\item \textbf{Beyond conics.}
Which projective incidence theorems for rational cubic and higher-degree
curves admit certificate-laden formulations that stay within quantifier-free
induction and avoid general root-existence principles?
\end{enumerate}
